% ====================================================================
% CUSTOM COMMANDS & ENVIRONMENTS
% ====================================================================
% File ini berisi custom command dan environment untuk mempermudah
% penulisan dan membuat dokumen lebih konsisten dan DRY (Don't Repeat Yourself).
%
% Tambahkan custom command di sini untuk:
% - Pengurangan kode yang berulang
% - Konsistensi formatting
% - Kemudahan modifikasi global
% ====================================================================

% --- COMMAND UNTUK CHAPTER TANPA NOMOR (Front Matter) ---
% Penggunaan lebih mudah daripada menulis \chapter* setiap kali
% Contoh: \frontmatterchapter{PERSEMBAHAN}
\newcommand{\frontmatterchapter}[1]{
  \chapter*{#1}
  \addcontentsline{toc}{chapter}{#1}
}

% --- COMMAND UNTUK KEYWORDS/KATA KUNCI ---
% Untuk konsistensi formatting kata kunci di abstrak
% Contoh: \keywords{Kata Kunci Satu, Kata Kunci Dua, Kata Kunci Tiga}
\newcommand{\keywords}[1]{
  \vspace{1cm}
  \noindent \textbf{Kata Kunci:} \textit{#1}
}

% --- COMMAND UNTUK KEYWORDS DALAM BAHASA INGGRIS ---
% Contoh: \keywordsen{Keyword One, Keyword Two, Keyword Three}
\newcommand{\keywordsen}[1]{
  \vspace{1cm}
  \noindent \textbf{Keywords:} \textit{#1}
}

% --- ENVIRONMENT UNTUK CENTERLINE (SHORTCUT) ---
% Penggunaan lebih singkat dari \begin{center} ... \end{center}
% Contoh: \mycenterline{\textit{``quoted text''}}
\newcommand{\mycenterline}[1]{
  \begin{center}
    #1
  \end{center}
}

% --- COMMAND UNTUK AUTHOR DENGAN AFFILIATION ---
% Berguna jika diperlukan dalam cover atau halaman judul
% Contoh: \authorWithAffiliation{Nama Mahasiswa}{NIM}{Program Studi}
\newcommand{\authorWithAffiliation}[3]{
  {\normalsize\bfseries #1}\\
  {\normalsize #2}\\
  {\normalsize #3}
}

% --- COMMAND UNTUK QUOTE YANG DIPUSATKAN ---
% Untuk mengformat quote atau motto dengan mudah
% Contoh: \quoteline{``In life, you either win or learn...''}{Nelson Mandela}
\newcommand{\quoteline}[2]{
  \begin{center}
    \textit{#1}\\[0.5cm]
    (\small\textit{#2})
  \end{center}
}

% --- COMMAND UNTUK HORIZONTAL RULE DEKORATIF ---
% Untuk separator visual dalam dokumen
\newcommand{\decorativeHRule}{
  \vspace{0.5cm}
  \centerline{\hrulefill}
  \vspace{0.5cm}
}

% --- COMMAND UNTUK TEMPAT TANDA TANGAN (SIGNATURE BLOCK) ---
% Berguna untuk halaman persetujuan
% Usage: \signatureBlock{Nama}{Tanggal}{Jabatan}
\newcommand{\signatureBlock}[3]{
  \vspace{1.5cm}
  \noindent
  #2\\
  \vspace{3cm}
  \noindent\line(1,0){100}\\
  \noindent\small #1\\
  \noindent\small #3
}

% --- COMMAND UNTUK INFORMASI PENELITIAN ---
% Shortcut untuk menampilkan info pembimbing, penguji, dll
% Contoh: \advisorInfo{Nama Pembimbing}{Gelar}
\newcommand{\advisorInfo}[2]{
  \noindent\textbf{#1}\\
  {\normalsize #2}
}

% --- COMMAND UNTUK BARIS NAMA & NIP PENGUJI/PEMBIMBING ---
\newcommand{\supervisorrow}[4]{%
  % #1 = nomor
  % #2 = jabatan (Pembimbing Utama / Penguji Utama / dst)
  % #3 = nama
  % #4 = NIP
  
  \quad #1. #2 & \\[0.3em]
  
  \quad \quad Nama : #3 & \signatureLine \\[0.3em]
  \quad \quad NIP \; : #4 & \\[1em]
}

% --- ENVIRONMENT TABEL PEMBIMBING & PENGUJI ---
\newenvironment{approvalboard}{%
  \noindent
  \begin{tabular}{@{} p{0.7\textwidth} p{0.3\textwidth} @{}}
}{%
  \end{tabular}
}

% --- GARIS TANDA TANGAN ---
\newcommand{\signatureLine}{(\dotfill)}

% --- COMMAND UNTUK SUBSUBSECTION KHUSUS (a., b., dll. - Normal font, Menjorok Lurus Angka Akhir Subsection) ---
% Penggunaan: \subsubsubsection{Judul Subsubsection}
\newcommand{\subsubsubsection}[1]{%
  \par\vspace{10pt}%
  \refstepcounter{subsubsection}%
  \hspace*{1.2cm}{\normalfont\normalsize \alph{subsubsection}. #1}\par\nopagebreak\vspace{4pt}%
}

% ====================================================================
% CATATAN PENGGUNAAN:
% - Gunakan custom command ini secara konsisten di seluruh dokumen
% - Jika perlu mengubah format, hanya perlu edit di sini (DRY principle)
% - Tambahkan command baru sesuai kebutuhan dokumen Anda
% ====================================================================
